\section{Research question and the meaning of the experiment}
\subsection{The geological motivation}
A thin section is a slice of rock made sufficiently thin for transmitted-light microscopy. A single section contains many mineral grains. Under plane-polarised light (PPL), useful observations include colour, absorption, texture and boundaries. Under cross-polarised light (XPL), brightness and interference colours can change as the stage rotates. Extinction is the dark orientation of an anisotropic grain between crossed polarisers. Pleochroism is a change in apparent colour or absorption with orientation in PPL.

The motivation is therefore straightforward: a geologist uses changing appearances, so a model might benefit from them too. However, more photographs do not automatically provide more independent information. A model also needs reliable alignment, appropriate regions, useful features, enough independent specimens and trustworthy evaluation targets.

\subsection{Questions tested}
\begin{enumerate}
\item \textbf{Primary:} under the same regions, descriptors, splits and target policy, does a 12-position angle-aware model outperform a single-position model?
\item \textbf{Secondary:} does knowing the angles help compared with a model seeing the same 12 positions without angle information?
\item \textbf{Diagnostic:} how much do target availability, class support, training instability and preprocessing restrict interpretation?
\end{enumerate}
The operational task is five-family classification. It is not full mineral-species identification, automatic discovery of every grain, or an independently validated replacement for a petrographer.

\subsection{Three kinds of evidence}
This report uses saved GPU predictions and evaluation files for quantitative outcomes; the current source snapshot for implementation explanations; and local images and masks for visual demonstrations. The historical GPU checkout and complete resolved regularised-run configuration were not archived together. Consequently, the current source is an explanatory snapshot, not a certified byte-for-byte copy of the historical training environment. Local regenerated label tables can differ from archived training evidence; historical counts below are taken from the run evidence.

There is no fresh training in preparing this report. The 72-position, unaveraged experiment is a proposed improvement, not a completed experiment. This separation prevents an implementation idea from being presented as a result.

\subsection{Vocabulary and units of independence}
\begin{tabularx}{\linewidth}{@{}lX@{}}
\toprule Term & Meaning in this project\\\midrule
Specimen & A rock sample identifier; the unit kept separate across train/test splits.\\
Site & A particular field of view within a specimen, such as \code{granite_r1}.\\
Stage position & One angle of microscope-stage rotation.\\
Image & One mode, PPL or XPL, at one position. Two images form a pair.\\
Region & A connected area proposed by segmentation; it need not equal one physical grain.\\
Weak target & A provisional family or allowed family set inferred from rules and metadata.\\
Embedding & A numerical image description produced by a frozen pretrained neural encoder.\\
Fold / seed & A specimen split / an initialisation and sampling randomisation.\\\bottomrule
\end{tabularx}
\takeaway{Thousands of images of a few rocks do not equal thousands of independent rock samples. The held-out specimen, rather than the image or grain, is central to evaluation.}

\section{Dataset provenance and organisation}
\subsection{Collection from the public Virtual Microscope}
The starting material comes from the public Virtual Microscope collection associated with The Open University \cite{ukvm}. It consists of existing microscope photographs and descriptive specimen metadata. The project did not capture these images with a local microscope.

The retained collection files indicate an automated download workflow: discover specimen records, resolve their viewer resources, enumerate rotation sites and PPL/XPL frame URLs, save images in a consistent hierarchy, and record failures and metadata. Viewer resource information and specimen indexes support this reconstruction. The original downloader, referenced in older documentation as \code{tools/ukvm_download.py}, is absent from the reviewed source tree. Its exact retry, throttling and URL-parsing implementation therefore cannot be asserted from this snapshot. The saved files establish what was downloaded more strongly than they establish every detail of the scraping procedure.

The relevant hierarchy is:
\begin{quote}\small\ttfamily
data/raw/ukvm\_dataset/\newline
\hspace*{1em}specimens.csv; rotation\_fields.csv; failures.csv\newline
\hspace*{1em}specimen\_index.json\newline
\hspace*{1em}specimens/<specimen>/\newline
\hspace*{2em}meta.json; viewer resources\newline
\hspace*{2em}rot/<site>/ppl/000.jpg ... 071.jpg\newline
\hspace*{2em}rot/<site>/xpl/000.jpg ... 071.jpg
\end{quote}
Original metadata and attribution must remain attached to redistributed material. Public access does not imply unrestricted reuse; collection-specific and third-party rights still apply. The companion repository excludes raw imagery and full data arrays. Figure provenance and image-rights notes are provided separately.

\subsection{How many observations are available?}
At a complete site the nominal raw angles are $0^\circ,5^\circ,\ldots,355^\circ$. Thus there are 72 positions and $72\times2=144$ photographs. Each image in this collection is approximately $490\times490$ pixels. PPL and XPL are different lighting observations, not duplicate copies.

\begin{table}[htbp]\centering
\caption{Data reduction and available supervision. Counts describe different units and should not be divided indiscriminately.}
\begin{tabularx}{\linewidth}{@{}rX@{}}\toprule Count & Stage / interpretation\\\midrule
115 & Specimens in the local Phase 0 inventory.\\
490 & Raw rotation folders before duplicate removal.\\
244 & Unique sites after removing 246 repeated folders.\\
242 & Sites meeting basic completeness and image checks.\\
233 & Sites passing alignment quality control.\\
22,389 & Automatically proposed regions in the segmentation collection.\\
5,065 & Regions represented in the archived full embedding run.\\
4,956 & Regions with pooled out-of-fold predictions; 112 specimens.\\
581 & Singleton-labelled test regions, from 77 specimens.\\
323 & Partial-set test regions, from 29 specimens.\\
4,052 & Predicted test regions whose teacher abstains.\\\bottomrule
\end{tabularx}
\end{table}
The difference between 5,065 and 4,956 is 109 mixed/excluded regions (66 mixed and 43 excluded in the archived accounting). Of the 4,956 predicted regions, only 904 have singleton or partial-set supervision. A large prediction count therefore does not imply a large labelled benchmark.

\subsection{Phase 0: clean the collection before measuring it}
Phase 0 builds an explicit inventory. It checks required files, decodability, consistent dimensions and completeness. Duplicate-site detection uses selected reference files, including PPL frame 000, XPL frame 000 and PPL frame 037, within a specimen. This removes exact repeated sites without intentionally merging distinct specimens. A hash of representative frames is an engineering duplicate check, not proof that every near-duplicate has been found.

The resulting \code{sites.csv} contains canonical site identifiers and usability information; \code{specimens_clean.csv} normalises specimen records. Missing image/whole-slide resources and incomplete sites are documented in the report. Two unusable sites in the inventory are \code{glahm134751_r2} and \code{lun35_r5}. An adjacent-frame PPL exposure-change check flags a relative jump above approximately 15\%; this is a warning rather than an automatic mineral classification. One site, \code{watson389_r2}, receives this flag.

The metadata include rock type and reported constituent minerals. These describe a \emph{specimen}. They do not identify the mineral occupying an individual pixel or region. Monomineralic flags were retained for description, but the V1 policy does not use a rock title alone to assign every region that mineral.

\subsection{Spatial calibration}
Scale information was estimated from whole-slide scale metadata and the relationship between whole-slide and rotation views. In simplified form,
\begin{equation}
u_{\rm rot}=\frac{1000}{\text{whole-slide pixels per mm}}
\times d_{\rm saved}\times s_{\rm rot\to WSI},
\end{equation}
where $u_{\rm rot}$ is micrometres per rotation-image pixel, $d_{\rm saved}$ accounts for saved-image downsampling, and $s_{\rm rot\to WSI}$ is the measured registration scale. Ten calibration specimens yielded a median near $4.181\,\mu$m/pixel and a coefficient of variation near 0.9\%. A 490-pixel image width is therefore about 2.05 mm under this approximation. The 464-pixel valid-circle diameter is about 1.94 mm.

This approximate spatial scale helps interpret grain sizes. It does not calibrate brightness, exposure, thin-section thickness or birefringence. Those are separate measurements.

\fig{raw.png}{1}{Running example: raw photographs from \code{granite_r1}. The spatial structures rotate while optical appearance also changes. The source photographs originate from the Virtual Microscope collection; this layout is a project illustration.}{fig:raw}
\FloatBarrier
